\documentclass[11pt]{article}
\usepackage{amsmath,amsthm,amssymb}
\usepackage[colorlinks]{hyperref}
\usepackage{tikz, pgfplots}
\pgfplotsset{compat=1.17}

\begin{document}
\title{Quick answer key to Asymptotics}
\author{ChatGPT (with inputs by Arjun)}
\date{24 February}
\maketitle

Use the table of contents below to skip to a specific part
without seeing spoilers to the other parts.

I just used ChatGPT to write this one quickly.
ChatGPT can make mistakes, so if you spot anything that's wrong, flag me to ask.

\tableofcontents



\newpage


%%%%%%%%%%%%%%%%%%%%%%%%%%%%%%%%%%%%%%%%%%%%%%%%%%%%%%%%%%%%
\section{Solving Recurrences}

\subsection{1. $T(n) = T(\sqrt{n}) + O(n)$}

This recurrence is not in standard Master Theorem form because the subproblem size is $\sqrt{n}$ rather than $n/b$.

Let us perform a change of variables.

Set
\[
n = 2^m.
\]
Define
\[
S(m) = T(2^m).
\]

Then:
\[
T(\sqrt{n}) = T(2^{m/2}) = S(m/2).
\]

Thus the recurrence becomes
\[
S(m) = S(m/2) + O(2^m).
\]

Now this is in Master Theorem form:

\[
S(m) = aS(m/b) + f(m)
\]
with:
\[
a=1, \quad b=2, \quad f(m)=\Theta(2^m).
\]

Compute:
\[
m^{\log_b a} = m^{\log_2 1} = m^0 = 1.
\]

We compare $f(m)=\Theta(2^m)$ to $1$.

Clearly $2^m$ grows exponentially faster than any polynomial. Thus we are in Master Theorem Case 3.

Therefore:
\[
S(m) = \Theta(2^m).
\]

Returning to $n=2^m$, we conclude:

\[
T(n) = \Theta(n).
\]

\[
\boxed{T(n)=\Theta(n)}
\]

%%%%%%%%%%%%%%%%%%%%%%%%%%%%%%%%%%%%%%%%%%%%%%%%%%%%%%%%%%%%
\newpage

\subsection{2. $T(n)=6T(n/3)+O(n^2)$}

Now this is exactly in Master Theorem form:

\[
a=6,\quad b=3,\quad f(n)=\Theta(n^2).
\]

Compute:
\[
n^{\log_b a} = n^{\log_3 6}.
\]

Note:
\[
\log_3 6 = \frac{\ln 6}{\ln 3} > 1.
\]

Numerically,
\[
\log_3 6 \approx 1.63.
\]

So:
\[
n^{\log_3 6} = \Theta(n^{1.63}).
\]

Compare $f(n)=\Theta(n^2)$.

Since $2 > 1.63$, we are in Master Theorem Case 3.

Thus:
\[
T(n)=\Theta(n^2).
\]

\[
\boxed{T(n)=\Theta(n^2)}
\]

%%%%%%%%%%%%%%%%%%%%%%%%%%%%%%%%%%%%%%%%%%%%%%%%%%%%%%%%%%%%
\newpage

\subsection{3. $T(n)=4T(n/5)+O(n)$}

Here:
\[
a=4,\quad b=5,\quad f(n)=\Theta(n).
\]

Compute:
\[
n^{\log_b a}=n^{\log_5 4}.
\]

Since $4<5$, we have:
\[
\log_5 4 < 1.
\]

Thus:
\[
n^{\log_5 4}= \Theta(n^{c}) \quad \text{with } c<1.
\]

Now compare $f(n)=\Theta(n)$.

Since $n$ grows polynomially faster than $n^c$, we are again in Case 3.

Hence:
\[
T(n)=\Theta(n).
\]

\[
\boxed{T(n)=\Theta(n)}
\]

%%%%%%%%%%%%%%%%%%%%%%%%%%%%%%%%%%%%%%%%%%%%%%%%%%%%%%%%%%%%
\newpage

\subsection{4. $T(n)=3T(\sqrt{n})+\log n$}

Again substitute $n=2^m$ and define $S(m)=T(2^m)$.

Then:
\[
S(m)=3S(m/2)+m.
\]

Master Theorem form:

\[
a=3,\quad b=2,\quad f(m)=\Theta(m).
\]

Compute:
\[
m^{\log_2 3}.
\]

Since $\log_2 3>1$, this is polynomial with exponent about $1.585$.

Compare $f(m)=\Theta(m)$.

Since $1 < \log_2 3$, we are in Case 1.

Therefore:
\[
S(m)=\Theta(m^{\log_2 3}).
\]

Return to $n=2^m$:

\[
m=\log n.
\]

Thus:
\[
T(n)=\Theta((\log n)^{\log_2 3}).
\]

\[
\boxed{T(n)=\Theta((\log n)^{\log_2 3})}
\]
\newpage
%%%%%%%%%%%%%%%%%%%%%%%%%%%%%%%%%%%%%%%%%%%%%%%%%%%%%%%%%%%%
\section{Asymptotic Comparisons}

\subsection{1. $f(n)=n\log n$, $g(n)=10n\log(10n)$}

Observe:
\[
\log(10n)=\log 10 + \log n.
\]

Thus:
\[
g(n)=10n(\log n + \log 10).
\]

As $n\to\infty$, $\log 10$ is constant.

Hence:
\[
g(n)=\Theta(n\log n).
\]

Therefore:
\[
f(n)=\Theta(g(n)).
\]

%%%%%%%%%%%%%%%%%%%%%%%%%%%%%%%%%%%%%%%%%%%%%%%%%%%%%%%%%%%%
\newpage

\subsection{2. $f(n)=n^{1.01}$, $g(n)=n\log^2 n$}

Canceling an $n$ on both sides and considering $f^*(n) = n^{0.1}$ and $g^*(n) = \log^2 n$.

Making the usual substitution $n = 2^k$, $g^*(n) = k^2$ while $f^*(n) = 2^(k/10)$.

Thus, $g^*(n) = O(f^*(n))$

Thus, $g(n) = O(f(n))$.\\

\textbf{Note:} AI partially failed at solving as it basically claimed that the answer is the answer.

%%%%%%%%%%%%%%%%%%%%%%%%%%%%%%%%%%%%%%%%%%%%%%%%%%%%%%%%%%%%
\newpage

\subsection{3. $f(n)=(\log n)^{\log n}$, $g(n)=\frac{n}{\log n}$}

Take logarithms:

\[
\log f(n)=\log n \cdot \log\log n.
\]

While:
\[
\log g(n)=\log n - \log\log n.
\]

Since:
\[
\log n \log\log n \gg \log n,
\]

we conclude:
\[
f(n) \gg g(n).
\]

Thus:
\[
g(n)=O(f(n)).
\]

%%%%%%%%%%%%%%%%%%%%%%%%%%%%%%%%%%%%%%%%%%%%%%%%%%%%%%%%%%%%
\newpage

\subsection{4. $f(n)=\log(n!)$, $g(n)=n\log n$}

We bound:

Upper bound:
\[
n! \le n^n
\Rightarrow \log(n!) \le n\log n.
\]

Lower bound:
\begin{align*}
n! &\ge n \times (n-1) \times (n-2) \times \dots \times n/2\\
&\ge n/2 \times n/2 \times n/2 \times \dots \times n/2 \\
&\ge (n/2)^{n/2}\\
\implies n! &\ge (n/2)^{n/2}\\
\implies \log(n!) &\ge n/2 \log(n/2)\\
&= n/2 (\log n - \log 2)\\
&= n/2 \log n - n/2 \log(2)\\
\implies \log(n!) &\ge O(n \log n)
\end{align*}

Thus:
\[
\log(n!) = \Theta(n\log n).
\]

\textbf{Note:} AI messed this up and just claimed theorem. Although, more prompting got this answer.

%%%%%%%%%%%%%%%%%%%%%%%%%%%%%%%%%%%%%%%%%%%%%%%%%%%%%%%%%%%%
\newpage

\subsection{5. $f(n)=\binom{n}{k}$, $g(n)=n^k$}

We start with:
\[
\binom{n}{k}
= \frac{n!}{k!(n-k)!}
= \frac{n(n-1)(n-2)\dots(n-k+1)}{k!}.
\]

Upper Bound:

Each term in the numerator satisfies
\[
n-i \le n \quad \text{for } 0 \le i \le k-1.
\]

Thus,
\[
n(n-1)\dots(n-k+1)
\le n^k.
\]

Dividing by $k!$:
\[
\binom{n}{k}
\le \frac{n^k}{k!}.
\]

Since $k!$ is a constant (because $k$ is fixed),
\[
\binom{n}{k} = O(n^k).
\]

Lower Bound:

For sufficiently large $n > 2k$, we have
\[
n-i \ge \frac{n}{2}
\quad \text{for all } 0 \le i \le k-1.
\]

(This holds whenever $n \ge 2k$.)

Thus,
\[
n(n-1)\dots(n-k+1)
\ge \left(\frac{n}{2}\right)^k.
\]

Dividing by $k!$:
\[
\binom{n}{k}
\ge
\frac{1}{k!}\left(\frac{n}{2}\right)^k
=
\frac{1}{k!2^k} n^k.
\]

Since $\frac{1}{k!2^k}$ is a constant,
\[
\binom{n}{k} = \Omega(n^k).
\]

Thus, $\binom{n}{k} = \Theta(n^k)$.

\textbf{Note:} AI messed this up and just claimed theorem. Although, more prompting got this answer.


%%%%%%%%%%%%%%%%%%%%%%%%%%%%%%%%%%%%%%%%%%%%%%%%%%%%%%%%%%%%
\newpage

\subsection{6. $f(n)=n$, $g(n)=n^{1+\sin n}$}

Since $\sin n\in[-1,1]$,

\[
g(n)\in[n^0,n^2].
\]

Thus infinitely often:

\[
g(n)=n^2
\quad \text{and} \quad
g(n)=1.
\]

Therefore neither function bounds the other asymptotically.

\[
\boxed{\text{Neither } f=O(g) \text{ nor } g=O(f).}
\]
\newpage
%%%%%%%%%%%%%%%%%%%%%%%%%%%%%%%%%%%%%%%%%%%%%%%%%%%%%%%%%%%%
\section{Harder Recurrences}

Most of these had to be edited quite a bit as GPT failed at them quite a bit.

%%%%%%%%%%%%%%%%%%%%%%%%%%%%%%%%%%%%%%%%%%%%%%%%%%%%%%%%%%%%
\subsection{1. $T(n)=\sum_{i=1}^{n-1}T(i)+1$}

Consider $T(n) - T(n-1) = (\sum_{i=1}^{n-1}T(i)+1) - (\sum_{i=1}^{n-2}T(i)+1) = T(n-1)$

Thus, $T(n) = 2T(n-1)$.

Thus, $T(n) = 2^n T(1) \implies T(n) = O(n)$.

\textbf{Note:} This was not hard. Somehow, GPT failed at it and just used induction after saying "notice the pattern".


%%%%%%%%%%%%%%%%%%%%%%%%%%%%%%%%%%%%%%%%%%%%%%%%%%%%%%%%%%%%
\newpage

\subsection{2. $T(n)=T(7n/10)+T(n/5)+O(n)$}

As

\[
\frac{7}{10}+\frac{1}{5}=0.7+0.2=0.9<1.
\]

Thus, $T(n) = T(7n/10)+T(n/5)+O(n) = T(49n/100) + T(n/25) + O(9n/10) + O(n)$. The sum is geometrically decreasing, hence we make the guess: $T(n) \le 10 n = O(n)$.

(B) For $n =1$, we have $T(1) = 1 < 10$.

(S) Let $T(k) \le 10k$ for $k < n$. We will show $T(n) \le 10 n$.

\begin{align*}
    T(n) &= T(7n/10)+T(n/5)+ n\\
    &\le 10 \frac{7n}{10} + 10 \frac{n}{5} + n\\
    &= 10 n
\end{align*}

Hence, by induction, we are done!

%%%%%%%%%%%%%%%%%%%%%%%%%%%%%%%%%%%%%%%%%%%%%%%%%%%%%%%%%%%%
\newpage

\subsection{3. $T(n)=T(\log n)+n$}

Expand:

\[
T(n)=n+\log n+\log\log n+\dots
\]

This series is dominated by the first term.

Thus, we conclude by induction:

\[
\boxed{T(n)=\Theta(n)}
\]

%%%%%%%%%%%%%%%%%%%%%%%%%%%%%%%%%%%%%%%%%%%%%%%%%%%%%%%%%%%%
\newpage

\subsection{4. $T(n)=T(n-\sqrt{n})+1$}

Our basic question is that "how many times do we need to subtract $\sqrt{n}$ before the number reaches some constant."

So what do we do? One idea is to try out some values and then guess. The other idea is to consider a continuos extension.
\begin{align*}
\frac{d n}{d t} &= - \sqrt{n}\\
-\frac{d n}{\sqrt{n}} &= d t\\
\int -\frac{d n}{\sqrt{n}} &= \int d t\\
t &= O(\sqrt{n})
\end{align*}

And we now complete by induction!


%%%%%%%%%%%%%%%%%%%%%%%%%%%%%%%%%%%%%%%%%%%%%%%%%%%%%%%%%%%%
\newpage

\subsection{5. $T(n)=T(n-1)+n^k$}

Unroll:

\[
T(n)=\sum_{i=1}^n i^k.
\]

If you remember the forms for $k = 1, 2, 3$ that is $\frac{n (n+1)}{2}, \frac{n (n+1) (2n + 1)}{6}$ and $(\frac{n (n+1)}{2})^2$ respectively; you could reasonably guess that this extends to the general case $T(n)=\Theta(n^{k+1})$ and induct away.

The "formal, non-guess way" is we compare to integral:

\[
\int_1^n x^k dx = \Theta(n^{k+1}).
\]

Thus, we conclude by induction

\[
\boxed{T(n)=\Theta(n^{k+1})}
\]

%%%%%%%%%%%%%%%%%%%%%%%%%%%%%%%%%%%%%%%%%%%%%%%%%%%%%%%%%%%%

\end{document}